\section{Structural Factors in Uneven Domain Coverage}
\label{sec:coverage-imbalance}

The distribution of scientific-agent studies varies substantially across the eleven domains: the current corpus records 172 work--domain memberships in materials science and nine in agricultural, food, forestry, animal, and fishery sciences. Interpreting this disparity requires attention to the conditions under which research tasks can be specified, executed, and validated. Relevant factors include mechanistic knowledge, experimental tractability, the availability of intervention data, representation of scientific states, and standards of evidence. Building on the cases in Section~\ref{sec:domains}, this section examines the disciplinary conditions that may shape these differences. The available evidence supports comparison of these conditions, but does not establish their respective causal effects on publication volume.

\subsection{Mechanistic knowledge, experimental tractability, and evaluation criteria}

The maturity of a scientific task has at least three distinct dimensions: knowledge of the underlying mechanism and its domain of applicability; the capacity to implement reproducible interventions and measurements; and the availability of criteria for discriminating among hypotheses or substantiating a scientific claim. Progress along these dimensions can be uneven. Experimental optimization may be feasible under incomplete mechanistic knowledge, whereas a theoretically characterized system may remain difficult to investigate because of measurement limitations or poorly constrained experimental conditions.

Chemical screening and materials formulation illustrate the former condition. The mobile robotic chemist selects photocatalyst mixtures using measured hydrogen-production outcomes, and autonomous electrolyte research integrates candidate selection with laboratory measurements~\cite{burger2020,asd0487}. Both provide an operationally defined search space and measurable objectives without requiring a complete mechanistic account of candidate performance. In physical inference, by comparison, governing equations must be supplemented by appropriate boundary conditions, calibrated observations, and sufficient information to identify the parameters or mechanisms of interest.

These distinctions suggest that the distribution of agent research depends on the relationship between unresolved scientific decisions and the structure available to support them. Tasks with defined interventions and informative evaluation can accommodate adaptive selection or coordination. Fully specified routines may already be supported by conventional automation, while tasks with unresolved measurement or evaluation criteria require additional scientific groundwork. A monotonic association between disciplinary maturity and publication volume should therefore not be assumed. Experimental feedback is considered here as a condition shaping research feasibility, rather than as a mandatory inclusion criterion.

\subsection{Data adequacy and alignment of scientific representations}

Data adequacy depends on the inferential objective, the unit of replication, and the coverage of relevant interventions and contexts. Repeated images from one plot, cells from a limited set of donors, or densely sampled observational records can increase measurement precision while leaving independent replication and intervention coverage largely unchanged. In clinical and environmental studies, causal identification additionally depends on assumptions about treatment or exposure assignment and potential confounding. Raw data volume is consequently an incomplete indicator of the evidence available for scientific decisions.

Decision-making also requires information that static datasets may not preserve. Property tables and published results can omit unsuccessful procedures, alternative actions, sample preparation, instrument state, or the rationale for experimental selection. These omissions restrict reconstruction of the relationship among scientific state, intervention, and outcome. Institutional access restrictions and incompatible recording practices can further limit the comparability and reuse of clinical, industrial, and field data. The availability of scientific records, their accessibility, and their suitability for evaluating agent decisions should therefore be assessed separately.

Multimodal complexity is similarly determined by the scientific relationships among representations. Chemical and materials investigations can combine molecular graphs, spectra, crystal structures, microscopy, processing histories, and instrument signals. Their integration requires correspondence at the specimen, temporal, and experimental-condition levels, together with preservation of physical quantities, units, and uncertainty. Protocol digitization illustrates how an explicit representation can connect descriptions to executable chemical operations~\cite{asd0600}. Specialized analytical tools can support the corresponding conversion from measurements to structured feedback. Excessive compression into textual summaries may remove information required for product identification or differential interpretation of a biological phenotype.

Compatibility with language-model representations can facilitate literature analysis, software interaction, and report preparation. Scientific execution additionally requires task-specific state estimation, constraints, and validation. These requirements apply across symbolic planning, probabilistic optimization, learned control, and LLM-based systems. Multimodal input capability alone provides limited assurance that observations are correctly aligned with experimental states or translated into admissible actions.

\subsection{Feedback informativeness, intervention constraints, and validation endpoints}

The scientific value of feedback depends on its capacity to reduce uncertainty about the research objective. Relevant attributes include measurement error, repeatability, evaluation latency, experimental throughput, and discrimination among competing explanations. Rapid simulation can support efficient exploration, subject to the adequacy of the model and objective. Physical investigations introduce additional constraints through irreversible interventions, destructive measurements, limited parallelization, and incomplete reproducibility of initial conditions. These constraints are particularly consequential in clinical research, ecological studies, marine deployment, and flight missions.

Evaluation endpoints also differ in their relationship to the intended scientific contribution. Electrolyte stability requires further evaluation before supporting claims about full-cell lifetime; molecular binding requires additional evidence for clinical benefit; phenotype extraction requires agronomic validation for claims about yield; and numerical convergence requires model validation for physical interpretation. Intermediate endpoints can support rigorous methodological studies while leaving subsequent scientific and translational questions unresolved. Differences in publication volume may therefore partly reflect the availability of tractable evaluation stages within each domain.

\subsection{Discipline-specific research conditions and limitations}

\paragraph{01: Formal, information, and computational sciences.}
Executable representations, reproducible software environments, and explicit evaluation procedures facilitate program search and computational experimentation. For program-search tasks such as FunSearch, feasibility depends on the availability of an evaluator that corresponds to the mathematical or computational claim~\cite{asd0512}. Performance under finite tests, formal correctness, and algorithmic novelty require distinct forms of evidence. Reusable execution infrastructure can support numerous workflow studies, while open mathematical discovery remains dependent on problem formulation, evaluator completeness, and independent verification.

\paragraph{02: Physics, astronomy, and cosmic sciences.}
Established models and instrument interfaces support calibration, parameter exploration, simulation, and observation scheduling. Instrument-facing investigations remain sensitive to drift, noise, approximation error, observability, and parameter identifiability, even when calibration actions can be automated~\cite{asd0090}. Astronomical investigations are additionally constrained by the predominantly observational nature of their research objects. Acquisition and parameter inference require distinct assessments of measurement quality, uncertainty, and model consistency~\cite{asd0136,asd0121}. The existence of governing equations therefore does not guarantee that the needed observation or parameter can be obtained reliably.

\paragraph{03: Chemical sciences.}
Standardized operations and analytical measurements support reaction screening and condition optimization under incomplete mechanistic knowledge. Digitized synthesis protocols formalize the experimental operations available to an agent~\cite{asd0600}. Transfer across experimental settings nevertheless depends on the representation of mixing, phase behavior, addition sequence, impurities, and apparatus. Reaction signals and yield measurements provide evidence at particular stages; product identity, purity, selectivity, and claims of a new transformation require appropriate characterization. The prevalence of tractable optimization tasks can support a substantial methodological literature, while broader chemical discovery remains constrained by experimental transferability and interpretation.

\paragraph{04: Materials science.}
Composition and formulation spaces often admit measurable objectives and repeated candidate testing, enabling adaptive optimization despite incomplete mechanistic understanding~\cite{burger2020,asd0487}. Generalization is constrained by the coupled relationships among composition, processing, microstructure, properties, and service conditions. Thermal history, defects, and interfaces can alter performance at a fixed composition. Computational admissibility and measured candidate performance leave further requirements for reproducible synthesis, devices, and service reliability~\cite{asd0001,burger2020}. The relatively high coverage of materials science is thus consistent with the availability of bounded optimization tasks, without implying comparable maturity across the full materials-development process.

\paragraph{05: Earth and environmental sciences.}
Earth-observation infrastructure makes geospatial analysis, event detection, and adaptive acquisition technically accessible~\cite{asd0650,asd0722}. Mechanistic interpretation is limited by heterogeneous spatial and temporal resolution, incomplete observations, sensor bias, and restricted intervention opportunities. Climate, land use, and human activity can jointly influence observed changes, complicating causal attribution. Controlled replication is limited at the system scale, while counterfactual simulation depends on model structure and parameterization. These conditions support a distinction between the feasibility of observational workflows and the evidence required for environmental explanations or transferable interventions.

\paragraph{06: Life sciences.}
Biological literature, molecular databases, and analysis software support hypothesis generation, single-cell analysis, and perturbation selection~\cite{asd0138,asd0020}. Mechanistic inference is constrained by dependence on cell type, genetic background, time, and environment, together with incomplete coverage of combinatorial perturbations. Large numbers of measured cells can coexist with limited independent biological replication. Existing perturbation datasets support retrospective selection studies, while prospective experiments test whether the selected intervention is informative under newly realized conditions~\cite{asd0020,king2004}. Data and software availability can sustain extensive analytical research while experimental confirmation and transfer across biological contexts remain demanding.

\paragraph{07: Medical and health sciences.}
Literature and electronic health records support analytical workflows, alongside imaging, physiological signals, laboratory measurements, and longitudinal clinical data. Treatment-selection mechanisms, disease severity, cohort selection, and follow-up processes can complicate causal inference from these records. Privacy restrictions and variation across institutions and populations further constrain external evaluation. Analytical and molecular endpoints permit intermediate evaluation, but clinical claims require evidence corresponding to the intended population, intervention, and outcome~\cite{asd0002,asd0054}. The availability of intermediate analytical and molecular endpoints can therefore coexist with substantial barriers to demonstrating safety, efficacy, and patient benefit.

\paragraph{08: Agricultural, food, forestry, animal, and fishery sciences.}
Food--molecule analysis and plant phenotyping provide defined tasks for tool use and computational investigation~\cite{asd0695,asd0510,asd0634}. Field-level inference is constrained by genotype--environment--management interactions and incomplete observation of biologically relevant states. Similar phenotypes can arise from pathogens, nutrient limitation, or water stress; imagery may provide insufficient information about root conditions or prior management. Independent plots, locations, and seasons remain necessary for evaluating transfer beyond the development setting. Forest growth, fish-population dynamics, and food processing introduce additional temporal and contextual dependencies. Controlled environments can improve experimental tractability, while field and production claims require broader validation. These constraints offer plausible explanations for low coverage, alongside limitations of literature collection.

\paragraph{09: Engineering and industrial technology sciences.}
Numerical solvers support simulation orchestration, and fabrication platforms permit experimentally evaluated design search~\cite{asd0047,asd0504}. Scientific validity depends on problem specification, including boundary conditions, constitutive assumptions, contact models, discretization, and operating regimes. Numerical convergence addresses only part of this assessment. Manufacturing introduces tolerances, defects, scale effects, and reliability requirements, while industrial data access can restrict replication. Established automation also supports many fully specified routines. Agent research may consequently concentrate on uncertain design, coordination across tools, and failure recovery, with publication counts reflecting these eligible research tasks rather than the overall extent of industrial automation.

\paragraph{10: Aerospace, marine, and transportation sciences.}
Communication limitations, constrained resources, and time-sensitive observation opportunities motivate autonomous planning and execution. Operational mission systems demonstrate that this demand can justify integrated analysis, planning, and execution~\cite{asd0645}. Such systems require decisions under partial observability using images, telemetry, position, energy, timing, and environmental information. Platform scarcity, consequential failures, and limited opportunities for deployment restrict exploratory experimentation and independent replication. Underwater communication and localization, and the transition from transportation simulation to operational testing, impose additional requirements. These conditions can produce a strong need for autonomy alongside comparatively slow development, validation, and publication cycles.

\paragraph{11: Social, behavioral, economic, and knowledge-system sciences.}
Scholarly records and linguistic interaction support empirical analysis and agent-based simulation~\cite{asd0624,asd0652,asd0786}. External validity depends on the empirical grounding of modeled preferences, incentives, and institutional arrangements. Multiple mechanisms may reproduce similar observations, and human adaptation can alter responses to interventions. Increasing prompt-defined simulated interaction expands evidence about the model under its assumptions; inference about human institutions still requires independent empirical evaluation~\cite{asd0786}. Linguistic accessibility therefore facilitates system construction while leaving behavioral identification and institutional transfer as substantive scientific limitations.

\begin{table}[H]
\caption{Research conditions and validation requirements across the eleven domains. Entries synthesize representative cases and disciplinary constraints; they do not constitute domain-wide measurements or causal estimates of publication volume.}
\label{tab:coverageimbalance}
\scriptsize
\begin{tabularx}{\textwidth}{@{}p{0.07\textwidth}>{\raggedright\arraybackslash}X>{\raggedright\arraybackslash}X>{\raggedright\arraybackslash}X@{}}
\toprule
Module & Tractable research tasks & Principal evidence constraints & Further validation requirements \\
\midrule
01 & Programs and computational experiments & Evaluator completeness and novelty & Independent verification or proof \\
02 & Calibration, simulation, observation & Observability and identifiability & Calibrated inference and model comparison \\
03 & Protocol execution and reaction screening & Operational transfer and product identity & Characterization and replication \\
04 & Formulation and property optimization & Processing history and scale dependence & Synthesis, devices, and service conditions \\
05 & Earth-observation analysis and planning & Data alignment and causal attribution & Regional transfer and model assessment \\
06 & Omics analysis and perturbation selection & Context dependence and intervention coverage & Prospective tests and biological replication \\
07 & Clinical analysis and candidate design & Treatment effects and population transfer & Evidence for specified clinical outcomes \\
08 & Food relationships and phenotyping & Latent states and independent field trials & Locations, seasons, and management regimes \\
09 & Simulation and experimental design search & Model specification and manufacturing variation & Physical validation and reliability testing \\
10 & Observation planning and resource allocation & Constrained exploration and scarce deployment & Validation under operational conditions \\
11 & Scholarly analysis and social simulation & Behavioral identification and external validity & Independent human or institutional evidence \\
\bottomrule
\end{tabularx}
\end{table}

\subsection{Cross-domain comparison and empirical assessment}

The proposed explanations should be examined within comparable research roles and validation settings. Materials formulation and field agriculture differ in the reproducibility of interventions and the transferability of local measurements to variable operating environments. Life science and medicine share methods and data resources, but biological assays and patient outcomes require distinct inferential and validation procedures. Aerospace further illustrates how operational demand for autonomy can diverge from opportunities for inexpensive experimentation. Comparisons at this level preserve scientifically relevant differences within and across domains.

An empirical assessment would require consistent coding of independent experimental units, observational and interventional evidence, specimen and temporal alignment, feedback latency and error, admissible actions, human decisions, and external validation. Relevant analyses include associations between intervention coverage and prospective evaluation, error differences between structured analytical outputs and raw multimodal inputs, and retention of improvements when moving from proxy objectives to independently evaluated outcomes. Comparisons among planning, optimization, LLM, and hybrid systems additionally require comparable actions, resources, and evaluators. These variables have not been uniformly coded across the present corpus, precluding quantitative attribution of publication differences to the proposed factors.

\subsection{Literature provenance and classification effects}

The collection procedures and non-exclusive counting rules in Sections~\ref{sec:methods} and~\ref{sec:corpus} qualify the structural interpretation above. Historical mission systems and field applications may be described outside contemporary scientific-agent terminology, while citation-network expansion increases the visibility of closely connected studies. These collection effects can coexist with differences in experimental feasibility. Distinguishing them requires task-level sampling and comparison; the present cases do not identify their independent contributions to publication volume.
